

We exploit a randomized control trial and a dynamic structural model to analyze how preferential college admissions affect pre-college effort and longer-term educational outcomes in Chile. The policy (PACE) guaranteed selective college admission to disadvantaged students graduating in the top 15\% of their high school class. Using a dataset of 9,006 students combining administrative and survey data, we find PACE increased first-year enrollment in selective colleges by 3.0 percentage points (35\% of the control mean), an effect waning to 1.5 percentage points (30\%) by the fifth year. The policy reduced students' pre-college effort, likely due to biased beliefs regarding the returns to effort in college admission and persistence. Counterfactual simulations from the model show policymakers could mitigate these unintended disincentives while preserving PACE’s college attainment gains by correcting students' beliefs about effort's returns in college persistence. Our results demonstrate that students' perceptions can critically shape the impacts of preferential admission policies.


